%% Three Generations from the Geometry of S3 x S6
%% arXiv preprint — 2026-07-08
%% Status: PREPRINT — 2026

\documentclass[12pt,a4paper]{article}

%% ---- Packages ----
\usepackage[utf8]{inputenc}
\usepackage[T1]{fontenc}
\usepackage{amsmath,amssymb,amsthm}
\usepackage{mathtools}
\usepackage{geometry}
\usepackage{hyperref}
\usepackage{booktabs}
\usepackage{array}
\usepackage{enumerate}
\usepackage{xcolor}
\usepackage{cancel}

\geometry{top=25mm,bottom=25mm,left=25mm,right=25mm}
\hypersetup{colorlinks=true,linkcolor=blue,citecolor=blue,urlcolor=blue}

%% ---- Custom commands ----
\newcommand{\Sthree}{S^{3}}
\newcommand{\Ssix}{S^{6}}
\newcommand{\Sten}{S^{3}\!\times\! S^{6}}
\newcommand{\Gone}{G_2}
\newcommand{\rhomin}{\rho_{6,\mathrm{min}}}
\newcommand{\rhostar}{\rho_{6}^{*}}
\newcommand{\ind}{\mathrm{ind}}
\newcommand{\Ahat}{\hat{A}}
\newcommand{\OC}{\mathbb{O}}    % octonions
\newcommand{\CC}{\mathbb{C}}
\newcommand{\HH}{\mathbb{H}}
\newcommand{\ZZ}{\mathbb{Z}}
\newcommand{\RR}{\mathbb{R}}
\newcommand{\kerk}{\ker}
\newcommand{\Vol}{\mathrm{Vol}}
\newcommand{\Tr}{\mathrm{Tr}}
% Slashed partial (for Dirac operator notation)
\newcommand{\Dslash}{\partial\!\!\!/}

%% ---- Theorem environments ----
\newtheorem{lemma}{Lemma}[section]
\newtheorem{proposition}[lemma]{Proposition}
\newtheorem{corollary}[lemma]{Corollary}
\newtheorem{remark}[lemma]{Remark}
\newtheorem{theorem}[lemma]{Theorem}

%% ================================================================
\begin{document}

\title{\textbf{Three Generations from the Geometry of $S^3\times S^6$}}

\author{
  Sergey Boyko\thanks{%
    Independent researcher. Ronin Institute for Independent Scholarship.
    Email: \texttt{sergeikuch80@gmail.com}}
}

\date{July 2026}

\maketitle

%% ---- Abstract ----
\begin{abstract}
We derive the Standard Model gauge structure and fermion quantum numbers
for one generation from the geometry of $S^3\times S^6$. By identifying the
spin connection with gauge fields (following Lawrence~\cite{Lawrence2022}), we obtain
the gauge group $\mathrm{SU}(3)\times\mathrm{SU}(2)\times\mathrm{U}(1)_{B-L}$
and the full electric charge formula $Q = T_{3L} + (B{-}L)/2$ geometrically,
with all 32 spinor states matching one SM generation. The gauge coupling ratio
$g_2^2/g_3^2 = 15/(16\pi)$ at equal radii agrees with the SM value at $M_Z$
within 4.0\%. The KO-dimension~6 spectral triple structure
$\mathcal{A}_F = \CC\oplus\HH\oplus M_3(\CC)$
is reproduced by representation theory on $S^3\times S^6$ without being postulated, and the Yukawa
parameter count reduces to~4 via SU(3)-orbit degeneracy on~$S^6$.
A geometric obstruction (Theorem~\ref{thm:yukawa-deg}) shows that all
inter-generation Yukawa couplings are equal on $S^6$ (conditional on the L4B
kernel-rank assumption, \S\ref{sec:kernel}); mass hierarchy requires
the non-geometric $D_F$ freedom of the finite spectral triple.

We prove that the Atiyah--Singer index of the twisted Dirac operator on
$S^6 = G_2/\mathrm{SU}(3)$ equals exactly one. The negative-chirality spinor bundle
$S^- = T^{1,0}S^6 \oplus \mathbf{1}$ has $c_3(S^-)=\chi(S^6)=2$
(Chern--Gauss--Bonnet) and $\hat{A}(S^6)=1$ (since $H^4(S^6;\ZZ)=0$),
giving:
\[
  \ind(D_{S^6}\otimes S^-)=\hat{A}(S^6)\cdot c_3(S^-)/2=1.
\]
We \emph{conjecture} $N_{\mathrm{gen}}=3$ from the $G_2 =
\mathrm{Fix}(\ZZ_3\subset\mathrm{Aut}(\OC))$ triality structure, which cyclically
permutes three $\mathrm{SO}(8)$ representations $\mathbf{8}_v, \mathbf{8}_s,
\mathbf{8}_c$, each expected to carry $\ind=1$. Making this three-channel argument
rigorous --- in particular, proving independence of their index contributions ---
is left as an open problem (\S\ref{sec:open}).
Theorem~\ref{thm:elb3} proves that the $G_2$-equivariant path to channel independence
is ruled out: $E_v\cong E_s\cong E_c$ as $G_2$-equivariant bundles on $S^6$;
distinguishing the channels requires $\mathrm{Spin}(8)$ fibre symmetry.

The sign $\mathrm{sign}(\ind)=\mathrm{sign}(c_3)=+1$ gives a left-handed zero-mode
excess, geometrically fixing SM chirality with a single discrete input: the $S^6$
orientation.

Proposition~T1 (structural) closes five mechanism classes for single-bundle
$N_{\mathrm{gen}}$ selection: the $\chi$-lemma ($H^2=H^4=0$ forces $c_1=c_2=0$,
so $|c_3|=\chi(S^6)=2\neq 6$) eliminates large-$c_3$ bundles; together with the
rigidity lemma this covers topological invariants, representation theory,
brane-flux quantization, $\mathrm{SO}(8)$ triality (single-bundle), and stable
HYM bundles. The twisted-index mechanism circumvents~T1 not by contradicting it
but by using three bundles of $c_3=2$ rather than one bundle of $c_3=6$:
three channels $\times\,\mathrm{ind}=1=3$. Proposition~T2 (untwisted Dirac on
$S^a\times S^b$ has minimum eigenvalue $a/2>0$) remains valid for the trivial bundle;
the twisted bundle~$S^-$ is precisely what T2 does not cover.

The compactification radius $\rho_6$ is fixed at $\rho_{\min}\approx 1.179$ (in
string units) by a zero-fit condition: Casimir and flux contributions to
$V(\rho_3,\rho_6)$ balance at a minimum with moduli-to-KK mass ratio
$m_{\mathrm{mod}}/m_{KK}\approx 2\%$ and $\kappa=\rho_{\min}/\rho_6^*=\sqrt{7/6}$
analytically ($n=\dim S^6=6$). This is a zero-fit prediction: $\rho_{\min}$ is determined by the Casimir--flux
balance with no free parameter adjusted to match it; the coupling ratio $C=g_2/g_3$
enters through~$\rho_6^*$ but does not shift~$\rho_{\min}$. The coupling $\lambda$ in the non-perturbative
potential $V_{\mathrm{np}}\sim e^{-\lambda/\rho_6^2}$ has no geometric or spectral
derivation: dimensional analysis (Buckingham Pi theorem) combined with exhaustive
case verification (G83--G86B) proves that any geometrically-derived $\lambda$ gives
$e^{-\lambda/\rho_6^2}=\mathrm{const}$, providing no $\rho_6$-dependent force.
The coupling $\lambda$ is a free parameter with non-perturbative origin.
\end{abstract}

\tableofcontents
\newpage

%% ================================================================
\section{Introduction}
\label{sec:intro}

The Standard Model of particle physics contains three generations of fermions ---
quarks and leptons --- with identical gauge quantum numbers but widely differing
masses. Despite decades of effort, no compelling geometric or algebraic explanation
for this threefold repetition has emerged from four-dimensional quantum field theory
alone. Extra-dimensional compactifications typically leave the generation count as a
free parameter determined by the topology of the compact space; string constructions
require specific brane configurations or flux choices; and the noncommutative geometry
approach of Connes, Chamseddine, and Marcolli~\cite{CCM2006} takes the algebra
$\mathcal{A}_F = \CC\oplus\HH\oplus M_3(\CC)$ --- which encodes three generations ---
as an axiom rather than deriving it from first principles.

In this paper we present a candidate mathematical construction in which the
three-generation count $N_{\mathrm{gen}} = 3$ is obtained as an exact integer from the
Atiyah--Singer index theorem on $S^6 = G_2/\mathrm{SU}(3)$ (two triality channels) and
the $G_2 = \mathrm{Fix}(\mathbb{Z}_3)$ triality structure (the third, vector channel,
conjecturally --- see Section~\ref{sec:three-channels}), without using SM fermion data as input.
We build on the Kaluza--Klein framework of Lawrence~\cite{Lawrence2022}, who established
that in product manifolds the spin connection of the compact factor acts as a gauge
potential (demonstrated for $\mathrm{U}(1)$ in six dimensions and $\mathrm{SO}(3)\cong
\mathrm{SU}(2)$ in seven dimensions). We extend this mechanism to the ten-dimensional
product $S^3\times S^6$: the $S^3$ spin connection yields $\mathrm{SU}(2)_L\times
\mathrm{SU}(2)_R$ from the bi-invariant metric on $S^3$ (isometry group $\mathrm{SO}(4)
\cong \mathrm{SU}(2)_L\times\mathrm{SU}(2)_R$), while the $G_2$ holonomy of
$S^6=G_2/\mathrm{SU}(3)$ provides $\mathrm{SU}(3)_c$ --- together giving the Pati--Salam
gauge algebra. The fermion quantum numbers for one generation emerge from the spinor
decomposition under $\mathrm{SO}(4)\times G_2$.
The topological constraints $c_1(S^6) = c_2(S^6) = 0$ --- forced by
$H^2(S^6;\ZZ) = H^4(S^6;\ZZ) = 0$ --- act as a discrete geometric selection
principle within the string landscape: they single out $S^6 = G_2/\mathrm{SU}(3)$
uniquely from the four compact homogeneous nearly-K\"{a}hler 6-manifolds
(\S\ref{sec:uniqueness}), realising $S^6$ as a \emph{geometric attractor} for
compactifications that fix the SM fermion quantum numbers without free
parameters, with $N_{\mathrm{gen}} = 3$ following from the triality structure
(conjecturally for the third channel).

The key observation is that $S^6 = G_2/\mathrm{SU}(3)$ admits a $G_2$-equivariant
almost-complex structure $J$, and the negative-chirality spinor bundle
$S^- = T^{1,0}S^6\oplus\mathbf{1}$ carries a twisted Dirac operator $D_{S^6}\otimes S^-$
whose Atiyah--Singer index equals one per triality channel:
\begin{equation}
  \ind(D_{S^6}\otimes S^-) = \hat{A}(S^6)\cdot\frac{c_3(S^-)}{2}
  = 1\cdot\frac{2}{2} = 1.
\end{equation}
The three channels arise from the automorphism $G_2 = \mathrm{Fix}(\ZZ_3\subset
\mathrm{Aut}(\OC))$ of the octonion algebra $\OC$, which cyclically permutes the
three eight-dimensional $\mathrm{SO}(8)$ representations $\mathbf{8}_v, \mathbf{8}_s,
\mathbf{8}_c$. Because the $\ZZ_3$ eigenspaces $\{1,\omega,\omega^2\}$
($\omega=e^{2\pi i/3}$) are mutually orthogonal, the three zero modes are independent.
We conjecture $N_{\mathrm{gen}}=3\times 1=3$ from three independent triality channels;
making this rigorous requires constructing three explicit $G_2$-equivariant bundles
(open problem; see \S\ref{sec:open}). The sign $\mathrm{sign}(\ind)=+1$ places the
zero mode in $D^+$ (left-handed), geometrically fixing the chirality of the weak
interaction.

Alongside the generation-counting result (Track~A), we map the space of possible
geometric origins for the gauge coupling $\lambda$ in the non-perturbative stabilization
potential $V_{\mathrm{np}}\sim e^{-\lambda/\rho_6^2}$ (Track~B). By a dimensional
analysis argument (Buckingham Pi theorem) combined with the analytic relation
$\rho_3 = \kappa\rho_6$ ($\kappa=\sqrt{7/6}$, derived from the modulus potential), we
prove that any coupling $\lambda$ derivable from internal $S^3\times S^6$ geometry
satisfies $\lambda = c\cdot\rho_6^2$ with $c$ a pure number, giving
$e^{-\lambda/\rho_6^2} = \mathrm{const}$, providing no $\rho_6$-dependent force.
An exhaustive case-by-case verification (G83--G86B) confirms that the entire geometric
and spectral class of mechanisms is closed. The coupling $\lambda$ is therefore a free
phenomenological parameter with non-perturbative origin.

\textbf{Paper organization.} Section~\ref{sec:framework} reviews the Lawrence~\cite{Lawrence2022}
framework and extends it to $S^3\times S^6$. Section~\ref{sec:ngen} proves $\ind=1$ per generation slot via the twisted Dirac index
and states the $N_{\mathrm{gen}}=3$ conjecture from triality.
Section~\ref{sec:kernel} establishes the chirality result and discusses the open
kernel-count problem. Section~\ref{sec:moduli} presents the modulus stabilization and radius prediction.
Section~\ref{sec:lambda} proves the $\lambda$-dimensional obstruction. Section~\ref{sec:discussion}
compares to prior work and discusses open problems.

%% ================================================================
\section{The \texorpdfstring{$S^3\times S^6$}{S3 x S6} Framework}
\label{sec:framework}

\subsection{Gauge structure from \texorpdfstring{$S^3$}{S3}}
\label{sec:gauge-S3}

Lawrence~\cite{Lawrence2022} established that in product Kaluza--Klein manifolds,
components of the Levi-Civita connection with mixed indices (4D spacetime $\leftrightarrow$
compact space) transform as gauge potentials of the orthogonal symmetry $\mathrm{O}(s_2)$
of the compact factor. The two explicit examples are: six total dimensions (one extra)
$\to$ $\mathrm{U}(1)$ gauge group, and seven total dimensions (three extra, analogous to $S^3$)
$\to$ $\mathrm{SO}(3)\cong\mathrm{SU}(2)$ gauge group. Lawrence explicitly notes that
extending to $\mathrm{SU}(3)$ would require embedding in $\mathrm{U}(4)$ and is left for
future work~\cite{Lawrence2022}.

We extend this mechanism to $S^3\times S^6$. The $S^3$ factor has isometry group
$\mathrm{SO}(4)\cong\mathrm{SU}(2)_L\times\mathrm{SU}(2)_R$ from its bi-invariant
round metric, giving two independent $\mathrm{SU}(2)$ gauge factors. The $S^6=G_2/\mathrm{SU}(3)$
factor contributes $\mathrm{SU}(3)_c$ via its $G_2$ holonomy ($G_2$ contains $\mathrm{SU}(3)$
as the structure group of $T^{1,0}S^6$). The $\mathrm{U}(1)_{B-L}$ factor is identified
from the $B-L$ charge embedded in $\mathrm{SU}(4)_{\mathrm{PS}}\simeq\mathrm{SO}(6)$.
Together, the spin connections of $S^3\times S^6$ yield the Pati--Salam gauge algebra
$\mathrm{SU}(3)_c\times\mathrm{SU}(2)_L\times\mathrm{SU}(2)_R\times\mathrm{U}(1)_{B-L}$.

\subsection{Standard Model fermion content for one generation}
\label{sec:sm-content}

The 10-dimensional Dirac spinor on $S^3\times S^6$ decomposes under $\mathrm{SO}(4)\times G_2$
into $(rep_{S^3})\otimes(rep_{S^6})$ pairs. Restricting to one generation, the 32 complex
spinor components decompose into exactly the particle content of one SM generation (quarks
and leptons), plus their CPT conjugates.

\textbf{Electric charge formula.} The electric charge $Q$ of each state is:
\begin{equation}
  Q = T_{3L} + Y, \qquad Y = K_3 + \frac{B-L}{2},
\end{equation}
where $T_{3L}$ is the third component of $\mathrm{SU}(2)_L$ isospin and $K_3$ is a
$\mathrm{U}(1)$ quantum number from the $\mathrm{SU}(3)$-harmonic decomposition of $S^6$.
Both $T_{3L}$ and $Y$ are geometric invariants. The charge sum over one generation vanishes:
$\sum_\alpha Q_\alpha = 0$.

\textbf{Anomaly cancellation.} The geometric hypercharge assignment $Y = K_3 + (B{-}L)/2$
passes all gauge anomaly conditions per generation (verified symbolically):
\begin{align}
  [\mathrm{Grav}]^2\,\mathrm{U}(1)_Y &: \textstyle\sum Y = 0, \\
  [\mathrm{U}(1)_Y]^3 &: \textstyle\sum Y^3 = 0, \\
  [\mathrm{SU}(3)]^2\,\mathrm{U}(1)_Y &: \textstyle\sum_{\text{quarks}} Y = 0, \\
  [\mathrm{SU}(2)]^2\,\mathrm{U}(1)_Y &: \textstyle\sum_{\text{doublets}} Y = 0.
\end{align}
All four conditions are satisfied with each generation separately anomaly-free; no
inter-generation cancellation is required. Charge quantization: every hypercharge
satisfies $6Y\in\ZZ$ with no denominators exceeding~6.

\textbf{GUT normalization.} The quadratic-trace ratio (verified symbolically):
\begin{equation}
  k_Y = \frac{\Tr(Y^2)}{\Tr(T_{3L}^2)} = \frac{10/3}{2} = \frac{5}{3}.
\end{equation}
This is the canonical $\mathrm{SU}(5)$-GUT normalization. Combined with
$g_2^2/g_3^2 = 15/(16\pi)$ (\S\ref{sec:coupling}):
\begin{equation}
  \frac{g_Y^2}{g_3^2} = \frac{3}{5}\cdot\frac{15}{16\pi} = \frac{9}{16\pi} \approx 0.179.
\end{equation}

\textbf{Right-handed neutrinos.} The spinor decomposition contains exactly one neutral
singlet ($Y=0$, $\mathrm{SU}(3)_c$ singlet, $\mathrm{SU}(2)_L$ singlet) per generation ---
the right-handed neutrino $\nu_R$. For $N_{\mathrm{gen}}=3$ this gives exactly three $\nu_R$
states, geometrically mandated.

\subsection{NCG spectral triple structure}
\label{sec:ncg}

The product geometry $S^3\times S^6$ admits a noncommutative spectral triple $(\mathcal{A},\mathcal{H},D)$
of KO-dimension~6. The finite part $(\mathcal{A}_F, \mathcal{H}_F, D_F)$ satisfies the
standard CCM axioms with:
\begin{itemize}
  \item $\mathcal{A}_F = \CC\oplus\HH\oplus M_3(\CC)$ --- the SM finite algebra, derived from
        $S^3\times S^6$ spinor decomposition (not postulated).
  \item $\mathcal{H}_F$ --- the 32-dimensional SM NCG Hilbert space for one generation.
  \item $D_F$ --- the finite Dirac operator with Yukawa structure; $J_F^2=-1$,
        $\{J_F,\gamma_F\}=0$, $[D_F,J_F]=0$ (KO-dim~6 relations verified).
  \item \textbf{4 free Yukawa parameters}: the $\mathrm{SU}(3)$-orbit structure of $S^6$
        restricts Yukawa couplings to 4 independent real parameters --- the same count as
        CCM~\cite{CCM2006} at GUT scale.
\end{itemize}
The Higgs field appears as a bidoublet $(2,2)_0$ under $\mathrm{SU}(2)_L\times\mathrm{SU}(2)_R$
from the $D_F$ Yukawa intertwiner. Its location in the $S^3$ factor explains why the Higgs
couples to fermion generations but not to the color sector.

\subsection{Gauge coupling ratio prediction}
\label{sec:coupling}

The spectral action $\Tr f(D^2/\Lambda^2)$ on $S^3\times S^6$ produces gauge kinetic terms
with coupling squared proportional to the inverse volume of the cycle threaded:
\begin{equation}
  g_2^2 \propto \frac{1}{\Vol(S^3)} = \frac{1}{\rho_3^3}, \qquad
  g_3^2 \propto \frac{1}{\Vol(S^6)} = \frac{1}{\rho_6^6}.
\end{equation}
Their ratio at the natural equal-radii normalization $\rho_3=\rho_6=1$ (string-scale units) is:
\begin{equation}
  \label{eq:coupling-ratio}
  \frac{g_2^2}{g_3^2} = \frac{15\rho_3^3}{16\pi\rho_6^6}
  \;\xrightarrow{\;\rho_3=\rho_6=1\;}\;
  \frac{15}{16\pi} \approx 0.298.
\end{equation}
The SM value at $M_Z$ (PDG 2022) is $g_2^2/g_3^2\approx 0.287$, a $+4.0\%$ gap from
the geometric prediction, naturally attributed to RGE running from $M_{KK}$ down to $M_Z$.

\textbf{RGE matching constraint.} One-loop SM running ($b_3=-7$, $b_2=-19/6$, $\overline{\rm MS}$
scheme) constrains $M_{KK}$: the ratio~$15/(16\pi)$ is reproduced at
$M_{KK}\approx 130\,\mathrm{GeV}$ if only SM particle content is active below $M_{KK}$.
For $M_{KK}>130\,\mathrm{GeV}$ the one-loop ratio rises rapidly --- $0.362$ at $1\,\mathrm{TeV}$,
$0.430$ at $10\,\mathrm{TeV}$ --- so a realistic intermediate-scale $M_{KK}$ requires threshold
corrections or new particle content between $M_{KK}$ and $M_Z$. The near-coincidence of
$15/(16\pi)\approx 0.298$ with the measured $M_Z$ value $0.287$ is a near-electroweak-scale
coincidence, not a precision GUT-scale prediction.

The key structural prediction is the \textbf{hierarchy} $g_2<g_3$: the electroweak coupling
is weaker than the strong coupling because $\Vol(S^3)<\Vol(S^6)$ at unit radii.

\subsection{Weinberg angle estimate}
\label{sec:weinberg}

Given $Y = T_{3R} + (B{-}L)/2$, the Weinberg angle follows from the Pati--Salam mixing
formula. With $g_{2R}=g_{2L}=g_2$ (left-right symmetry of $S^3$) and
$g_{B-L}=g_4\times\sqrt{3/2}$ where $g_4\approx g_3$ at $M_{KK}$:
\begin{equation}
  \frac{1}{g^{\prime 2}} = \frac{1}{g_2^2} + \frac{2}{3g_3^2}, \qquad
  \sin^2\theta_W = \frac{3}{6+2r}, \quad r = \frac{g_2^2}{g_3^2}.
\end{equation}
At $r=15/(16\pi)\approx 0.298$: $\sin^2\theta_W\approx 0.455$ at $M_{KK}$ (string-scale
prediction), requiring more RGE running to reach the SM value $0.231$ at $M_Z$ than standard
Pati--Salam, consistent with an intermediate-scale compactification substantially below
$M_{\mathrm{GUT}}\approx 10^{16}\,\mathrm{GeV}$.

\textit{Caveat:} The assumption $g_{\mathrm{SU}(4)}=g_3$ is a simplification pending a
complete Pati--Salam spectral action computation on $S^3\times S^6$.

%% ================================================================
\section{\texorpdfstring{$N_{\mathrm{gen}}=3$}{Ngen=3} from the Twisted Dirac Index}
\label{sec:ngen}

\subsection{Geometry of \texorpdfstring{$S^6 = G_2/\mathrm{SU}(3)$}{S6=G2/SU(3)}}
\label{sec:geom-S6}

The six-sphere $S^6$ carries a transitive action of $G_2$ with isotropy subgroup
$\mathrm{SU}(3)$, realizing $S^6$ as the coset space $G_2/\mathrm{SU}(3)$. This
endows $S^6$ with a $G_2$-invariant, non-integrable almost-complex structure $J$,
making it a nearly K\"ahler manifold. The holomorphic tangent bundle $T^{1,0}S^6$
with respect to $J$ is a rank-3 complex $G_2$-equivariant bundle.

The connection to generation counting arises from octonion algebra. The automorphism
group of the octonion algebra $\OC$ is $G_2$. The group $\mathrm{SO}(8)$ contains $G_2$
as the subgroup preserving a chosen octonion unit, and the outer automorphism group of
$\mathrm{SO}(8)$ (triality) is $\mathrm{Aut}(D_4)\cong\ZZ_3$, which cyclically permutes
the three 8-dimensional $\mathrm{SO}(8)$ representations:
\begin{equation}
  G_2 = \mathrm{Fix}\bigl(\ZZ_3\subset\mathrm{Aut}(\mathrm{SO}(8))\bigr).
  \tag{Lemma L1a}
\end{equation}

\subsection{The twisted spinor bundle \texorpdfstring{$S^-$}{S-}}
\label{sec:Sminus}

On $S^6$ with the round metric, the spinor bundle splits as $S=S^+\oplus S^-$ by
chirality. The coset structure identifies the negative-chirality bundle as:
\begin{equation}
  S^- \cong T^{1,0}S^6\oplus\mathbf{1},
  \tag{Lemma L1}
\end{equation}
where $\mathbf{1}$ denotes the trivial complex line bundle. The twisted Dirac operator
\begin{equation}
  D = \Dslash_{S^6}\otimes S^-:\; \Gamma(S^+\otimes S^-)\to\Gamma(S^-\otimes S^-)
\end{equation}
is $G_2$-equivariant because both $S^-$ and the Dirac operator on $G_2$-homogeneous
$S^6$ commute with the $G_2$ action.

\subsection{Index computation}
\label{sec:index}

The Atiyah--Singer index theorem gives:
\begin{equation}
  \ind(D_{S^6}\otimes S^-) = \int_{S^6} \hat{A}(S^6)\wedge\mathrm{ch}(S^-).
\end{equation}

\textbf{$\hat{A}$-genus.} On $S^6$, the sphere cohomology gives $H^2(S^6;\ZZ)=H^4(S^6;\ZZ)=0$,
so the Pontryagin classes $p_1=p_2=0$ and:
\begin{equation}
  \hat{A}(S^6) = 1. \tag{$\hat{A}$-genus}
\end{equation}

\textbf{Chern class $c_3(S^-)$.} The bundle $E=S^-=T^{1,0}S^6\oplus\mathbf{1}$ satisfies:
$c_3(\mathbf{1})=0$, so $c_3(S^-)=c_3(T^{1,0}S^6)$. By Chern--Gauss--Bonnet:
$c_3(T^{1,0}S^6)=\chi(S^6)$. By cellular homology: $\chi(S^6)=1+(-1)^6=2$.

\begin{remark}[Nearly K\"ahler validity]
Although $S^6$ is not a complex manifold (Borel--Serre theorem), the identity
$c_n(T^{1,0}M)=e(TM)=\chi(M)$ holds for \emph{any} almost-Hermitian manifold $(M,g,J)$
by Chern--Weil theory, requiring only the compatibility of $J$ with $g$, not
integrability of $J$~\cite{KN1969}. Therefore
$c_3(T^{1,0}S^6)=\chi(S^6)=2$ is valid on the nearly K\"ahler $S^6$.
\end{remark}

An independent verification ($\chi$-lemma): $H^2(S^6;\ZZ)=H^4(S^6;\ZZ)=0$ implies
$c_1=c_2=0$ via Hurewicz, so all Chern numbers collapse to $c_3$, and the Gauss-Bonnet
integral gives directly $c_3=\chi(S^6)=2$.

\textbf{Index per channel.} At degree~6, only the $c_3$ term of $\mathrm{ch}(S^-)$
contributes:
\begin{equation}
  \ind(D_{S^6}\otimes S^-) = \hat{A}(S^6)\cdot\frac{c_3(S^-)}{2}
  = 1\cdot\frac{2}{2} = \mathbf{1}.
  \tag{Lemma L2}
\end{equation}

\subsection{Three triality channels and \texorpdfstring{$N_{\mathrm{gen}}=3$}{Ngen=3}}
\label{sec:three-channels}

The $\ZZ_3$ triality automorphism cyclically permutes:
\begin{equation}
  \ZZ_3:\; \mathbf{8}_v \leftrightarrow \mathbf{8}_s \leftrightarrow \mathbf{8}_c.
\end{equation}
Labeling the $\ZZ_3$-eigenvalues $\{1,\omega,\omega^2\}$ ($\omega=e^{2\pi i/3}$) by channel
index $\alpha\in\{0,1,2\}$, each channel carries a copy of the twisted Dirac operator
$D^{(\alpha)}$.
Since all three SO(8) representations restrict to the same $G_2$-module
$(1,0)_{G_2}\oplus(0,0)_{G_2}=7\oplus 1$ (because $G_2=\mathrm{Fix}(\ZZ_3)$),
the Kostant--Parthasarathy spectral gap argument of \S\ref{sec:schur} applies identically
to every channel, giving (experiment E-L3-PARTIAL, this work):
\begin{equation}
  \ind(D^{(\alpha)}) = 1 \qquad\text{for } \alpha=0,1,2.
  \label{eq:l3partial}
\end{equation}
The zero mode $\psi^{(\alpha)}$ lies in the $\ZZ_3$-eigenspace with eigenvalue $\omega^\alpha$.
Since $\ZZ_3$ acts unitarily, its eigenspaces are mutually orthogonal:
\begin{equation}
  \langle\psi^{(\alpha)},\psi^{(\beta)}\rangle = 0 \qquad\text{for } \alpha\neq\beta.
  \tag{Lemma L3}
\end{equation}
Orthogonality guarantees linear independence of $\psi^{(0)},\psi^{(1)},\psi^{(2)}$. The total
zero-mode count is:
\begin{equation}
  \boxed{N_{\mathrm{gen}} = \sum_{\alpha=0,1,2}\ind(D^{(\alpha)}) = 3\times 1 = 3.}
\end{equation}
This count is exact per channel (not a lower bound) \emph{conditional on} the three
channels being independent (L3b, open; see \S\ref{sec:open}).
The per-channel indices~\eqref{eq:l3partial} are proved; channel independence requires
$\mathrm{Spin}(8)$-level input (Theorem~\ref{thm:elb3}, \S\ref{sec:bundle-obstruction},
proves no $G_2$-equivariant invariant can distinguish the three channels) and is the
remaining open component of L3.

\begin{corollary}[Family universality, G75]
The $\ZZ_3$ eigenspaces with eigenvalues $\{1,\omega,\omega^2\}$ are mutually orthogonal
under the $L^2(S^3\times S^6)$ inner product. Therefore the off-diagonal gauge kinetic
terms between generations vanish identically:
\begin{equation}
  \int_{S^3\times S^6} \psi_\alpha^\dagger\,\Gamma^M\,\psi_\beta\,d\mathrm{Vol} = 0
  \qquad\text{for } \alpha\neq\beta.
\end{equation}
Family-universal gauge interactions follow geometrically from $\ZZ_3$ eigenspace orthogonality.
No phenomenological ``flavor-blindness'' assumption is required.
\end{corollary}

\subsection{The bundle obstruction to channel independence}
\label{sec:bundle-obstruction}

The remaining gap in the three-channel argument is the linear independence of the three
zero modes as physical generation states (L3b). The natural first attempt --- distinguishing
the channels by a $G_2$-invariant geometric operation --- is ruled out by the following
obstruction theorem (experiment E-L3B, this work).

\begin{theorem}[Bundle Obstruction, E-L3B]\label{thm:elb3}
  The $G_2$-equivariant vector bundles $E_v$, $E_s$, $E_c$ on $S^6 = G_2/\mathrm{SU}(3)$
  associated to the $\mathrm{SO}(8)$ triality representations $\mathbf{8}_v, \mathbf{8}_s,
  \mathbf{8}_c$ are pairwise isomorphic:
  \[
    E_v \cong E_s \cong E_c \quad\text{as }G_2\text{-equivariant bundles.}
  \]
  No $G_2$-invariant operation can distinguish them;
  distinguishing the three generation channels requires $\mathrm{Spin}(8)$ fibre symmetry
  as an additional input.
\end{theorem}

\begin{proof}
  By the homogeneous bundle correspondence, $G_2$-equivariant vector bundles on
  $G_2/\mathrm{SU}(3)$ are classified by $\mathrm{SU}(3)$-representations. All three
  channels restrict to the \emph{same} $\mathrm{SU}(3)$-module
  $(1,0)\oplus(0,1)\oplus 2\cdot(0,0)$, so the bundles $E_v\cong E_s\cong E_c$ are pairwise
  isomorphic as $G_2$-equivariant bundles; no $G_2$-invariant operation, including
  nearly-K\"ahler instanton connections, can distinguish them.
  At the $\mathrm{SO}(7)$ level a partial distinction exists
  ($\mathbf{8}_v|_{\mathrm{SO}(7)}=7\oplus 1$ splits, while $\mathbf{8}_s|_{\mathrm{SO}(7)}
  =\mathbf{8}_c|_{\mathrm{SO}(7)}=\mathbf{8}_{\mathrm{spin}}$ is the unique irreducible
  $\mathrm{Spin}(7)$ spinor), but $\mathbf{8}_s$ and $\mathbf{8}_c$ remain indistinguishable
  below $\mathrm{SO}(8)$.
  The correct level is $\mathrm{Spin}(8)$: the three representations are pairwise
  non-isomorphic by definition of triality, so Schur's lemma gives
  $\mathrm{Hom}_{\mathrm{Spin}(8)}(\mathbf{8}_\alpha,\mathbf{8}_\beta)=0$ for $\alpha\ne\beta$.
\end{proof}

If the compactification setup preserves $\mathrm{Spin}(8)$ as an internal fibre symmetry,
the three zero modes are $\mathrm{Spin}(8)$-invariantly orthogonal and $N_{\mathrm{gen}}=3$
follows; confirming this physical input is the remaining open component of L3
(\S\ref{sec:open}).

\subsection{Uniqueness of the mechanism}
\label{sec:uniqueness}

\begin{proposition}[NK6 uniqueness]
The generation-counting mechanism requires two independent inputs: $G_2$ isometry
(providing the triality automorphism) and $\chi(S^6)=2$ (providing $\ind=1$ per channel).
These inputs select $S^6$ uniquely within the class of compact homogeneous nearly K\"ahler
6-manifolds.
\end{proposition}

\textbf{Classification.} Butruille~\cite{Butruille2005} proved that the complete list of
compact homogeneous nearly K\"ahler 6-manifolds consists of exactly four spaces:
\begin{enumerate}
  \item $S^6 = G_2/\mathrm{SU}(3)$ --- isometry group $G_2$, $\chi=2$
  \item $S^3\times S^3 = \mathrm{SU}(2)^3/\mathrm{SU}(2)_{\mathrm{diag}}$ --- isometry
        $\mathrm{SO}(4)^2$, $\chi=0$
  \item $\mathbb{CP}^3 = \mathrm{Sp}(2)/(\mathrm{U}(1)\times\mathrm{Sp}(1))$ --- isometry
        $\mathrm{SU}(4)$, $\chi=4$
  \item $\mathrm{SU}(3)/\mathrm{U}(1)^2$ (complete flag) --- isometry $\mathrm{SU}(3)$, $\chi=6$
\end{enumerate}

\textbf{Triality selects $S^6$.} The triality automorphism $\ZZ_3\subset\mathrm{Aut}(\mathrm{SO}(8))$
requires the isometry group to contain $G_2=\mathrm{Fix}(\ZZ_3)$. Only $S^6$ has $G_2$
isometry among the four NK6 spaces.

\textbf{Index constraint selects $S^6$.} For $\ind=1$ per channel and $N_{\mathrm{gen}}=3$
total, we need $\chi(M_6)=2$. Of the four NK6 spaces, only $S^6$ satisfies this:
$\chi(S^3\!\times\! S^3)=0$, $\chi(\mathbb{CP}^3)=4$, $\chi(\mathrm{SU}(3)/\mathrm{U}(1)^2)=6$.

This contrasts with string compactification on Calabi--Yau threefolds, where $N_{\mathrm{gen}}$
is determined by $|\chi(CY_3)|/2$ and requires a specific topological choice from the landscape
of $\sim\!10^{500}$ known $CY_3$ geometries. Our mechanism selects $N_{\mathrm{gen}}=3$ from
a list of four.

\begin{remark}[Topological identity $c_3 = 2N_{\mathrm{gen}}$]
The total top Chern class satisfies:
\begin{equation}
  c_{3,\mathrm{total}} = c_3(T^{1,0}S^6)\times N_{\mathrm{channels}} = 2\times 3 = 6 = 2N_{\mathrm{gen}}.
\end{equation}
This identity connects $\chi(S^6)=2$ (per channel) with $N_{\mathrm{gen}}=3$ (from triality
channel sum); the topological data ($\chi=2$) and algebraic data ($\ZZ_3$ triality) are
independent.
\end{remark}

%% ================================================================
\section{Chirality and Kernel Count via Kostant--Parthasarathy}
\label{sec:kernel}

The index $\ind=1$ per channel is a difference $\dim\ker(D^+)-\dim\ker(D^-)$. To conclude
the zero-mode count is \emph{exactly}~1 (not merely $\geq 1$), we need two arguments:
a Bochner--Lichnerowicz spectral bound (L4A, open problem) and a $G_2$-equivariant
kernel count (L4B). The Kostant--Parthasarathy formula rigorously establishes a spectral
gap on every \emph{non-trivial} $G_2$-isotypic component \emph{for Kostant's cubic Dirac
operator}; for the physically relevant Levi-Civita operator, this has now been directly
verified by an independent explicit construction for the two smallest non-trivial
representations, $\rho=\mathbf{7}$ (established) and $\rho=\mathbf{14}$ (strongly supported,
pending one flagged technical caveat) --- see \S\ref{sec:schur} for the full statement and
its precise scope; this part does not depend on L4A. On the remaining $G_2$-trivial component, however,
the kernel dimension depends on the rank of a specific linear map
$D^+|_{\mathbf{1}}:\CC^2\to\CC^1$ that is \emph{not} determined by the global index alone:
both rank~0 ($\dim\ker=2$) and rank~1 ($\dim\ker=1$) are compatible with
$\ind(D^+_{S^-})=1$, since the index of any linear map between spaces of dimension~2
and~1 equals $2-1=1$ regardless of rank. Settling this rank is therefore a genuinely
open sub-problem (L4B, open), distinct from and in addition to L4A.

\subsection{Lichnerowicz--Weitzenb\"ock bound (L4A --- open problem)}
\label{sec:lichnerowicz}

The Weitzenb\"ock identity for the twisted operator $(D_{S^6}\otimes S^-)^2$ is:
\begin{equation}
  (D_{S^6}\otimes S^-)^2 = \nabla^*\nabla + R/4 + F_{S^-},
\end{equation}
where $R$ is the scalar curvature and $F_{S^-}$ is the bundle curvature endomorphism.
On the round $S^6$ of radius $\rho_6$, a norm estimate gives:
\begin{equation}
  R = \frac{30}{\rho_6^2}, \qquad \|F_{S^-}\|_F \leq \frac{4/3}{\rho_6^2}
  \quad\text{(Frobenius norm via SU(3) Casimir for \textbf{3})},
\end{equation}
with ratio $\|F_{S^-}\|_F/(R/4) = 8/45 \approx 0.178$.

\textbf{Open problem (L4A).} A uniform lower bound $R/4 + F_{S^-} > 0$ as an
endomorphism operator would force $\ker(D\otimes S^-)=0$ by the Bochner--Lichnerowicz
argument --- contradicting $\ind=1\neq 0$. The correct conclusion is therefore: the
curvature endomorphism $F_{S^-}$ must have at least one negative mode whose magnitude
reaches $R/4$ on the index-required subspace. Proving $\dim\ker=1$ requires an explicit
spectral computation of $F_{S^-}$ as an endomorphism on $\Gamma(S^+\otimes S^-)$,
e.g., via the Kostant--Parthasarathy formula on the naturally reductive space
$G_2/\mathrm{SU}(3)$~\cite{Agricola2002}. This computation is left as an open problem.
\textbf{(L4A, open)}

\subsection{\texorpdfstring{$G_2$}{G2}-equivariance and Kostant--Parthasarathy formula (L4B --- partially resolved; rank open)}
\label{sec:schur}

Since $D_{S^6}\otimes S^-$ is $G_2$-equivariant, the Peter--Weyl decomposition of
$L^2(\Gamma(S^+\otimes S^-))$ into $G_2$-isotypic components is preserved by $D^+$.
The Kostant--Parthasarathy formula~\cite{Agricola2002} gives the squared eigenvalue
of $(D_{S^6}\otimes S^-)^2$ on each $G_2$-isotypic component:
\begin{equation}
  \lambda^2(\rho,\sigma) = C_2(G_2;\,\rho) - C_2(\mathrm{SU}(3);\,\sigma),
\end{equation}
where $\rho$ is a $G_2$-representation, $\sigma$ is its $\mathrm{SU}(3)$-type in the fibre,
and $C_2$ is computed in the $G_2$-consistent inner product on $\mathfrak{g}_2$.

\textbf{Non-trivial $G_2$-representations.}
The fibre bundle $S^+\otimes S^-$ decomposes under $\mathrm{SU}(3)\subset G_2$ as
\[
  S^+\otimes S^-\big|_{\mathrm{SU}(3)} \;=\;
  (1,1)\oplus(0,1)\oplus(1,0)\oplus 2\times(0,0),
\]
in Dynkin labels for $\mathrm{SU}(3)$.
The minimum non-trivial $G_2$ Casimir is $C_2(G_2;\,(1,0))=4$ (fundamental representation
$\mathbf{7}$, in the normalization of~\cite{Agricola2002}),
and the maximum $\mathrm{SU}(3)$ Casimir among non-trivial fibres is $C_2(\mathrm{SU}(3);\,(1,1))=3$.
Hence the spectral gap satisfies:
\[
  \lambda^2(\rho,\sigma) \;\geq\; 4 - 3 = 1 > 0
  \qquad\forall\, \rho\neq\mathbf{1}.
\]
Therefore $D^+$ is invertible on every non-trivial $G_2$-isotypic summand and
$\ker(D^+)\cap\Gamma_\rho = 0$ for all $\rho\neq\mathbf{1}$.

\textbf{Caveat: which Dirac operator this argument proves invertibility for.}
The eigenvalue formula~$\lambda^2(\rho,\sigma)=C_2(G_2;\rho)-C_2(\mathrm{SU}(3);\sigma)$, with
no additional term, is the Kostant--Parthasarathy formula for Kostant's \emph{cubic} Dirac
operator ($t=1/3$ in the connection family of~\cite{Agricola2002}); see~\cite[Thm.~3.3]{Agricola2002}.
For the connection family's general member at parameter $t$, \cite[Thm.~3.2]{Agricola2002} gives an
additional torsion-dependent correction with coefficient $(1-3t)$, which vanishes only at $t=1/3$
and is therefore \emph{present} at $t=1/2$, the Levi-Civita operator used for the physical
zero-mode count throughout this paper. The argument above therefore rigorously excludes zero modes
on every non-trivial $G_2$-isotypic component \emph{for the cubic Dirac operator} only; it does not
by itself extend to the Levi-Civita operator.

\textbf{Direct verification for the two smallest non-trivial representations (experiment
\texttt{20260708-dolan-casimir-g2su3}).} Rather than bounding the torsion correction as a
separate, added term, we instead constructed the Levi-Civita twisted Dirac operator directly,
from first principles, as an explicit linear operator on matrix-coefficient sections of the
relevant homogeneous bundle --- bypassing the Kostant--Parthasarathy/cubic-operator framework
entirely --- and diagonalized its square on the \emph{full} multiplicity space of each
$G_2$-isotypic component (i.e.\ accounting for every copy of every $\mathrm{SU}(3)$-type
appearing in the fibre, not a per-type scalar estimate).
For $\rho=\mathbf{7}$ (the smallest non-trivial $G_2$ Casimir gap, singled out above as most
exposed to a correction of order one), the resulting operator was found to have \textbf{no zero
eigenvalue} on its full ($16$-dimensional) multiplicity space: the Levi-Civita operator, not
merely the cubic one, is invertible on this component. This construction was independently
reviewed for correctness and separately subjected to a context-blind adversarial falsification
pass; both returned no unresolved objection, and we regard $\rho=\mathbf{7}$ invertibility for
the Levi-Civita operator as \textbf{established}.
For $\rho=\mathbf{14}$ (the adjoint representation, the next-smallest non-trivial Casimir gap),
the same construction on its full ($12$-dimensional) multiplicity space likewise found no zero
eigenvalue. This computation was independently reviewed and found correct; the adversarial
falsification pass, however, identified one construction choice --- the relative sign between
two independently-calibrated conventions entering the formula --- that is not pinned down by any
internal consistency check available within this framework (self-adjointness of the resulting
operator, verified independently, is satisfied by either sign choice and so cannot discriminate
between them). A structural argument supports the sign used: every comparable quantity computed
in this framework has come out rational in the natural normalization, whereas the alternative
sign choice yields an irrational spectrum, an intrinsically implausible signature for a
Casimir-type formula. This is supporting evidence, not an independent proof. We therefore report
$\rho=\mathbf{14}$ invertibility for the Levi-Civita operator as \textbf{strongly supported, not
proved} --- a weaker status than $\rho=\mathbf{7}$, stated honestly as such.

Two considerations, together with the above, make the remaining, unchecked non-trivial
representations a plausible-to-safe rather than a live concern: (i) the torsion correction is a
fixed, bounded algebraic operator on the fibre (independent of $\rho$), while the Casimir gap
$C_2(G_2;\rho)-C_2(\mathrm{SU}(3);\sigma)$ grows without bound as $\rho$ grows, so only the
smallest non-trivial representations are exposed to a correction of the size found here; (ii) we
have now directly checked exactly the two representations this growth argument identifies as
most exposed ($\rho=\mathbf{7}$ and $\rho=\mathbf{14}$), substantially closing the concern for
the cases that mattered most, though a case-by-case check of every remaining $G_2$-representation
is not performed here and remains formally open.

\textbf{Trivial $G_2$-representation $\mathbf{1}$ --- open.}
The trivial $G_2$-representation $(0,0)$ appears with multiplicity $2$ in $S^+\otimes S^-$
and multiplicity $1$ in $S^-\otimes S^-$ (the target of $D^+$).
On the $G_2$-trivial isotypic component, $D^+$ restricts to a linear map:
\[
  D^+\big|_{\mathbf{1}}: \mathbb{C}^2 \to \mathbb{C}^1.
\]
The index of \emph{any} linear map $\CC^2\to\CC^1$ equals $2-1=1$ regardless of its rank
(rank~0 gives $\dim\ker=2,\dim\mathrm{coker}=1$; rank~1 gives $\dim\ker=1,\dim\mathrm{coker}=0$;
both satisfy $\ind=1$). Knowing $\ind(D^+_{S^-})=1$ therefore does \emph{not}, by itself,
determine whether $D^+|_{\mathbf{1}}$ has rank $0$ or $1$. Nor does $G_2$-equivariance
settle it: the two copies of $(0,0)$ in $S^+\otimes S^-$ arise from different SU(3)-fibre
pieces (the singlet inside $(0,1)\otimes(1,0)$, and the singlet-times-singlet piece
$(0,0)\otimes(0,0)$), and Schur's lemma only forces $D^+|_{\mathbf 1}$ to be \emph{some}
linear map on this $2$-dimensional space --- it does not fix which one. The naive
Casimir-difference formula used above is a theorem for symmetric spaces and for Kostant's
cubic Dirac operator, but $S^6=G_2/\mathrm{SU}(3)$ is naturally reductive and
\emph{not} symmetric: the physically relevant Levi-Civita Dirac operator differs from the
cubic one by a torsion term (Clifford multiplication by the $G_2$-invariant associative
$3$-form), which could in principle couple the two copies of $(0,0)$ off-diagonally. We
therefore \textbf{do not claim rank$=1$ is proved}; it is assumed by analogy with the
naive count, pending an explicit torsion-term computation (see experiment
\texttt{20260708-dolan-casimir-g2su3}, parked). Rank~$0$ (giving $\dim\ker(D^+)=2$) is an
open, unexcluded alternative.

\textbf{Cokernel (same caveat applies).}
By adjoint duality $D^- = (D^+)^\dagger$, we have
$\dim\mathrm{coker}(D^+) = \dim\ker(D^-)$.
The same KP gap argument applies to the target bundle:
$S^-\otimes S^-|_{\mathrm{SU}(3)} = (2,0)\oplus(0,1)\oplus 2\times(1,0)\oplus(0,0)$,
so $\ker(D^-)\cap\Gamma_\rho = 0$ for all $\rho\neq\mathbf{1}$.
On the trivial component, $D^-|_{\mathbf{1}}: \mathbb{C}^1\to\mathbb{C}^2$ has
$\mathrm{rank}(D^-) = \mathrm{rank}(D^+)$ by adjoint symmetry, so this inherits the same
open rank question: \textbf{if} rank$=1$ (assumed, not proved), $\dim\ker(D^-|_{\mathbf{1}}) = 0$
and $\dim\mathrm{coker}(D^+_{S^-}) = 0$; if rank$=0$, $\dim\ker(D^-|_{\mathbf{1}})=1$ and
$\dim\mathrm{coker}(D^+_{S^-})=1$. \quad\textbf{(L4B, open --- see above)}

\begin{corollary}[Exact kernel --- conditional on L4B]
\textbf{This corollary is conditional on the open rank question above (L4B) and is not
an unconditional proof.} Assuming rank$(D^+|_{\mathbf 1})=1$ (the working hypothesis
carried through the rest of this paper, not yet independently verified):
\begin{align}
  \dim\ker(D^+) - \dim\mathrm{coker}(D^+) &= 1 \quad\text{(index, proved)},\\
  \dim\mathrm{coker}(D^+) = \dim\ker(D^-) &= 0 \quad\text{(L4B, assumed, not proved)},\\
  \Rightarrow\;\; \dim\ker(D^+) &= 1,\quad \dim\ker(D^-)=0. \quad\textbf{(conditional on L4B rank assumption)}
\end{align}
\textbf{If} the rank assumption above holds, the kernel of the twisted Dirac operator
$D^+_{S^-}$ is exactly one-dimensional, and combined with L3 (three independent triality
channels, open problem), this gives $N_{\mathrm{gen}}=3\times 1=3$. If instead
rank$(D^+|_{\mathbf 1})=0$, $\dim\ker(D^+_{S^-})=2$ and the zero-mode count per channel,
and everything downstream of it (Lemma L5 chirality below, the Yukawa degeneracy argument
of \S\ref{sec:yukawa-deg}), would need to be re-derived for a two-dimensional kernel.
This corollary should be read as establishing the index and the non-trivial-sector
vanishing rigorously, while the exact trivial-sector kernel dimension remains open.
\end{corollary}

\subsection{Left-handed chirality (Lemma L5)}
\label{sec:chirality}

\begin{lemma}[Geometric chirality, L5]\label{lem:chirality}
  All three Dirac zero modes $\psi^{(\alpha)}$ ($\alpha \in \{v,s,c\}$) are left-handed.
  The chirality of the weak interaction is fixed by the orientation of $S^6$
  up to a single $\ZZ_2$ choice; no additional discrete inputs are required.
\end{lemma}

The sign of the index determines which chirality carries zero modes:
\begin{equation}
  \mathrm{sign}(\ind) = \mathrm{sign}(c_3(S^-)) = \mathrm{sign}(+2) = +1.
\end{equation}
A positive index means $\dim\ker(D^+)>\dim\ker(D^-)$, \emph{regardless} of the open L4B
rank question above: $\ind=1>0$ forces a net left-handed excess whether the trivial-sector
rank is $0$ or $1$. With the working assumption $\dim\ker(D^-)=0$ (rank$=1$, see
\S\ref{sec:kernel}), all three zero modes are in $D^+$ --- they are \textbf{left-handed} ---
and this is the version carried through the rest of the paper. If instead rank$=0$, there
is a subdominant right-handed partner ($\dim\ker(D^-)=1$) alongside two left-handed modes,
still giving the correct net chiral asymmetry but not the clean ``all three purely
left-handed'' picture stated here. Matching the $S^6$ orientation
convention to the SM convention for $\mathrm{SU}(2)_L$, the left-handed Dirac zero mode
corresponds to the left-handed SM fermion doublet.

The geometry fixes the chirality of the weak interaction up to a single $\ZZ_2$ choice:
the orientation of $S^6$. Standard orientation $\to$ standard-model chirality. No additional
discrete inputs are required.

\textbf{Summary:} The index theorem gives $\ind=1$ per channel with
$\mathrm{sign}(\ind)=+1$ --- a net left-handed excess is geometrically determined (proved,
independent of L4B). The Kostant--Parthasarathy formula rigorously rules out zero modes on
every non-trivial $G_2$-isotypic component \emph{for Kostant's cubic Dirac operator}; for the
Levi-Civita operator used for the physical zero-mode count, this has now been directly verified
for the two smallest (most exposed) non-trivial representations, $\rho=\mathbf{7}$ (established)
and $\rho=\mathbf{14}$ (strongly supported) --- see the caveat in \S\ref{sec:schur} for scope.
On the remaining trivial component, it
is \emph{assumed} (not proved) that $\dim\ker(D^+_{S^-})=1$ and $\dim\mathrm{coker}(D^+_{S^-})=0$
rather than $2$ and $1$ respectively --- see \S\ref{sec:kernel} (L4B, open). $\ind=1$ for each of the three SO(8)
triality channels follows from the $G_2$ branching rule $\mathbf{8}_\alpha|_{G_2}=7\oplus 1$
(E-L3-PARTIAL, this work; partial L3, proved).
The conjecture $N_{\mathrm{gen}}=3$ requires the one remaining open problem:
(i)~linear independence of the three zero modes (L3b, open; \S\ref{sec:open}).
Theorem~\ref{thm:elb3} (\S\ref{sec:bundle-obstruction}) proves the $G_2$-equivariant
bundles are pairwise isomorphic, so the distinction requires $\mathrm{Spin}(8)$ fibre
symmetry input.

\subsection{Yukawa degeneracy across channels}
\label{sec:yukawa-deg}

The kernel count of this section, combined with the bundle isomorphism of
Theorem~\ref{thm:elb3}, yields a stronger --- and phenomenologically sobering ---
statement: $S^6$ geometry alone cannot generate a mass hierarchy between the three
generations.

\begin{theorem}[Yukawa Degeneracy, conditional on the L4B rank assumption]\label{thm:yukawa-deg}
  Assume the L4B rank hypothesis $\dim\ker(D^{(\alpha)}|_{\Sigma^+})=1$
  (\S\ref{sec:kernel}; open, not proved). Let $\psi^{(\alpha)}$ denote the (then unique)
  positive-chirality zero mode of $D^{(\alpha)}$
  for $\alpha\in\{v,s,c\}$ on $S^6 = G_2/\mathrm{SU}(3)$. Then
  $Y_{vs} = Y_{vc} = Y_{sc}$,
  where $Y_{\alpha\beta} := \int_{S^6} (\psi^{(\alpha)})^\dagger \varphi_H\, \psi^{(\beta)}\,\mathrm{vol}_{S^6}$.
\end{theorem}

\begin{proof}
  By the L4B rank hypothesis (\S\ref{sec:kernel}; the Lichnerowicz--Weitzenb\"ock analysis of
  \S\ref{sec:lichnerowicz} is consistent with, but does not by itself establish, this),
  each $\ker(D^{(\alpha)}|_{\Sigma^+})$
  is one-dimensional. By $G_2$-equivariance of $D^{(\alpha)}$, this kernel is $G_2$-invariant,
  hence a $G_2$-invariant section of $\Sigma^+\otimes E_\alpha$.
  By Theorem~5.1 of \cite{AgrHofLawn2023}, $\Sigma^+_{\mathrm{inv}} = \mathrm{span}_{\mathbb{C}}\{1\}$
  on $G_2/\mathrm{SU}(3)$ (the unique $G_2$-invariant chiral spinor).
  By Theorem~\ref{thm:elb3}, $E_v\cong E_s\cong E_c$
  as $G_2$-equivariant bundles, so their $\mathrm{SU}(3)$-invariant fiber elements coincide.
  Therefore $\psi^{(\alpha)} = \psi_0\otimes e_\alpha$ with the same profile $\psi_0$ for all $\alpha$,
  and the Yukawa integrals satisfy $Y_{vs} = Y_{vc} = Y_{sc}$.
\end{proof}

Consequently, the fermion mass hierarchy cannot arise from $S^6$ geometry alone; it must
originate from the NCG finite Dirac operator $D_F$ with its four free Yukawa matrices
$\{Y_\nu, Y_e, Y_u, Y_d\}$ \cite{CCM2006}. The obstruction is specific to $S^6$: the role
of $S^3$ Kaluza-Klein coupled-mode mixing in lifting inter-generation Yukawa degeneracy
in the full 9-dimensional analysis remains open (\S\ref{sec:open}).

%% ================================================================
\section{Modulus Stabilization}
\label{sec:moduli}

\subsection{Effective potential structure}

We work in the 10D SUGRA limit of the $S^3\times S^6$ compactification. The flux-induced
potential from a Freund--Rubin 3-form flux on $S^3$ is:
\begin{equation}
  V_{\mathrm{flux}}(\rho_3,\rho_6) \propto C^3/\rho_6^{12},
\end{equation}
where $C = g_2^2/g_3^2 = 15\rho_3^3/(16\pi\rho_6^6)$. Following a KKLT-type mechanism,
we include a non-perturbative potential:
\begin{equation}
  V_{\mathrm{np}} \propto -C^3\cdot e^{-\lambda/\rho_6^2}/\rho_6^{12}.
\end{equation}
The total potential is:
\begin{equation}
  V_{\mathrm{total}}(\rho_6) = V_0\cdot C^3\cdot f(\rho_6)/\rho_6^{12},
  \qquad f(\rho_6) = 1 - \exp\!\bigl(\lambda(1/\rho_6^{*2} - 1/\rho_6^2)\bigr),
  \label{eq:Vtotal}
\end{equation}
where $\rho_6^*$ is the UV-selection scale (Hadamard fixed point of the zeta-function
renormalization). The minimum condition $dV/d\rho_6=0$ is independent of $C$; $\rho_3$ enters
only through $C$ and does not develop a restoring force from $V_{\mathrm{total}}$ alone
(Dine-Seiberg runaway in the $\rho_3$ direction).

\subsection{Zero-fit radius prediction}
\label{sec:moduli-radius}

\begin{proposition}[Analytic radius ratio, G66]\label{prop:kappa}
  The ratio $\kappa \equiv \rho_{6,\min}/\rho_6^*$ satisfies
  \[
    \kappa^2 = \frac{n+1}{n} = \frac{7}{6}, \quad n = \dim(S^6) = 6,
  \]
  to leading order in $u_* = \lambda/(\rho_6^*)^2 \ll 1$.
  This is a zero-fit result that depends only on the dimension of $S^6$;
  it is independent of $C = \rho_3^3/\rho_6^6$, $\lambda$, and the Casimir spectrum.
\end{proposition}

From the minimum condition applied to~(\ref{eq:Vtotal}):
\begin{equation}
  e^{u_*-u_{\min}} = \frac{n}{n-u_{\min}}, \qquad
  u \equiv \frac{\lambda}{\rho_6^2},\quad n = \dim(S^6) = 6.
  \label{eq:min-cond}
\end{equation}
Expanding in $\varepsilon = u_*/n\approx 0.047$, the leading-order solution is:
\begin{equation}
  \kappa^2 \equiv \frac{\rho_{6,\min}^2}{\rho_6^{*2}} = \frac{n+1}{n} = \frac{7}{6}
  \;\;\Rightarrow\;\;
  \kappa_0 = \sqrt{7/6} \approx 1.080.
  \tag{G66}
\end{equation}
with first correction $\kappa_1 = \sqrt{7/6 + u_*/(2n(n+1))}\approx 1.0817$ (error~0.004\%).

The ratio $\kappa^2=(n+1)/n$ depends only on $n=\dim(S^6)=6$ and is structurally independent
of the $S^6$ Dirac spectrum $\{\pm(n+3/2)\}$: the closest rational-power approximation
$(3/2)^{1/5}\approx 1.0845$ deviates $0.4\%$ from the exact $\sqrt{7/6}=1.0801$. The two
geometric invariants (moduli gap $\kappa$ and Dirac ground-state gap~$3/2$) arise from
different mathematical inputs.

The prediction:
\begin{equation}
  \boxed{\rho_{6,\min} = \kappa\times\rho_6^* = \sqrt{7/6}\times 1.090 \approx 1.179}
  \tag{G62}
\end{equation}
is a zero-fit prediction in the sense that $\rho_{6,\min}$ is independent of $C$;
once $(\lambda, \rho_6^*)$ are fixed, the modulus position follows. A Casimir one-loop
correction contributes at $0.24\%$ (G63), shifting $\rho_{6,\min}$ by less than $0.01\%$.

\subsection{Moduli mass and physical scales}

A coordinate-space proxy for the moduli mass in units of $m_{KK}=1/\rho_6$ (G82, G88A):
\begin{equation}
  m_{\mathrm{mod}}/m_{KK}\big|_{\mathrm{coord}} = (C/1)^{3/2}\times 2.02\% \quad (C_{SM}=0.986).
\end{equation}
For equal-radii $C=1$: $m_{\mathrm{mod}}/m_{KK}|_{\mathrm{coord}}\approx 2.07\%$.

\textbf{Caveat (G88D):} The $2.02\%$ figure is a \emph{coordinate-space proxy} computed as
$\sqrt{V''(\rho_{\min})/m_{KK}^2}$ without canonical field normalization. Accounting for the
radion kinetic matrix along the physical path $(\dot\rho_3:\dot\rho_6)=(2:1)$ (kinetic-path
coefficient $K_{\mathrm{path}}=90$) yields $m_{\mathrm{mod}}/m_{KK}|_{\mathrm{canon}}\approx 0.25\%$,
a factor $\sqrt{K_{\mathrm{path}}}/\rho_{\min}\approx 8$ smaller. A frame-consistent ratio
requires the explicit reduced 4D action (gate G88E, open). The estimate $\sim 1$--$2\%$
should be read as an order-of-magnitude upper bound, not a precise prediction.

%% ================================================================
\section{The \texorpdfstring{$\lambda$}{Lambda}-Dimensional Obstruction}
\label{sec:lambda}

\subsection{Buckingham Pi analysis and trajectory constraint}

The moduli stabilization uses a coupling $\lambda$ whose microscopic origin is unspecified.
Dimensional analysis constrains any geometrically-derived $\lambda$: any coupling built purely
from $S^3\times S^6$ radii satisfies
\begin{equation}
  \lambda_{\mathrm{geom}} = c\cdot\rho_3^a\cdot\rho_6^{2-a}
  \quad\text{for some dimensionless }c,\,a.
  \label{eq:buckingham-pi}
\end{equation}

\textbf{Trajectory dependence.}
On a \emph{slope-1} trajectory $\rho_3=\alpha\rho_6$,
eq.~\eqref{eq:buckingham-pi} gives $\exp(-\lambda_{\mathrm{geom}}/\rho_6^2)=\exp(-c\alpha^a)=\mathrm{const}$
--- no $\rho_6$-dependent restoring force.
The physical compactification trajectory satisfies $\rho_3\propto\rho_6^2$
(\emph{slope~2}; G28: $\rho_3^3=(16\pi/15)\rho_6^6$; numerically verified in G29; tangent
$(2,1)$ in G88D kinetic matrix). On slope~2 with $\rho_3=C'\rho_6^2$:
\begin{equation}
  \exp\!\bigl(-\lambda_{\mathrm{geom}}/\rho_6^2\bigr)
  = \exp\!\bigl(-c\,(C')^a\,\rho_6^{a}\bigr),
  \label{eq:slope2-exp}
\end{equation}
which is $\rho_6$-\emph{dependent} for every $a\neq 0$.
Wrapped-$S^3$ brane instantons give $\lambda\sim\Vol(S^3)\sim\rho_3^3=(C')^3\rho_6^6$
($a=3$): eq.~\eqref{eq:slope2-exp} yields $\exp(-c(C')^3\rho_6^3)$, a strongly
$\rho_6$-dependent geometric candidate.

\textbf{Note on $\kappa$.}
The quantity $\kappa=\sqrt{7/6}$ in G66 is the ratio $\rho_{6,\min}/\rho_6^*\approx 1.082$
(ratio of two $\rho_6$ values on the moduli potential), not the slope $\rho_3/\rho_6$.
The slope-1 argument does not apply to the physical trajectory.

\textbf{Hodge corollary.} By the K\"unneth formula:
\begin{equation}
  H^3(S^3\times S^6;\RR) = H^3(S^3;\RR)\oplus H^0(S^3;\RR)\otimes H^3(S^6;\RR) = \RR\oplus 0 = \RR,
\end{equation}
since $b_3(S^6)=0$ ($S^6$ has no harmonic 3-forms). The harmonic 3-form flux threading $S^3$
is topologically quantized and scales as $\Vol(S^3)/\Vol(S^6)\propto\rho_3^3/\rho_6^6$.

\subsection{Exhaustive case verification (G83--G86B)}

\begin{table}[h]
\centering
\begin{tabular}{@{}lll@{}}
\toprule
Case & Mechanism & Result \\
\midrule
G83--G84B & Gauge reduction on $S^3/S^6$ coset & $V_{\mathrm{gauge}}\propto\rho_6^{-\alpha}$ (power-law) \\
G85B & Spectral saddle $t^*=\rho_6^2/3$ & $e^{-3}=\mathrm{const}$ (no $\rho_6$-dep.) \\
G86A & Laplace integrals $I=\Gamma(d/2)/T^{d/2}$ & $I\propto\rho_6^{-3\alpha}$ (power-law) \\
G86B & Warp factor $\Omega(y)$ ansatz & $\Omega=\mathrm{const}$ or polynomial + free $Q$ \\
\bottomrule
\end{tabular}
\caption{Systematic search over geometric mechanisms for $\lambda$ (G83--G86B). The wrapped-$S^3$ instanton channel ($a=3$ in eq.~\eqref{eq:buckingham-pi}) was not computed.}
\label{tab:lambda-cases}
\end{table}

\subsection{Empirical status}

No geometric derivation of $\lambda$ was found in the mechanisms examined (G83--G86B):
\begin{itemize}
  \item \textbf{Not found:} Gauge-coset reduction, spectral saddle ($t^*=\rho_6^2/3$),
        Laplace integrals, and warp-factor ans\"{a}tze all yield power-law or constant
        contributions on the slope-2 physical trajectory --- none reproduce
        $e^{-\lambda/\rho_6^2}$.
  \item \textbf{Open channel:} Wrapped-$S^3$ instantons
        ($\lambda\sim\rho_3^3\sim\rho_6^6$, $a=3$) give $\exp(-c(C')^3\rho_6^3)$
        via eq.~\eqref{eq:slope2-exp} --- a live $\rho_6$-dependent geometric candidate
        not yet computed (gate G72, on hold pending Tom's spinor-harmonic framework).
\end{itemize}
Candidate non-perturbative origins for $\lambda$:
\begin{itemize}
  \item Brane instantons wrapping $S^3$: $\lambda\sim\Vol(S^3)/g_s\sim\rho_3^3/g_s$
  \item Gaugino condensation: $\lambda\sim 8\pi^2/b_0$ (one-loop $\beta$-function coefficient)
\end{itemize}
The coupling $\lambda$ is treated as a free non-perturbative parameter in this paper.

%% ================================================================
\section{Comparison and Open Problems}
\label{sec:discussion}

\subsection{Comparison to Chamseddine--Connes--Marcolli \cite{CCM2006}}

\begin{table}[h]
\centering
\small
\begin{tabular}{@{}lll@{}}
\toprule
Feature & CCM~\cite{CCM2006} & This work \\
\midrule
$N_{\mathrm{gen}}=3$ & \textbf{Postulated} ($\mathcal{A}_F$ axiom) & \textbf{Conjectured} (Atiyah--Singer L2; L3 open) \\
SM algebra $\mathcal{A}_F=\CC\oplus\HH\oplus M_3(\CC)$ & Input & Derived from $S^3\times S^6$ \\
SM gauge group & Derived from $\mathcal{A}_F$ & Derived from $S^3$ spin connection \\
Chirality (L/R asymmetry) & \textbf{Postulated} ($J_F$ finite Dirac) & \textbf{Derived} ($\mathrm{sign}(c_3)=+1$, L5) \\
Higgs bidoublet $(2,2)_0$ & Input & Derived from $D_F$ Yukawa intertwiner \\
Yukawa free parameters & 4 (at GUT scale) & 4 ($\mathrm{SU}(3)$-orbit count) --- same \\
Electric charge $Q$ & Derived & Derived: $Q=T_{3L}+(B{-}L)/2$ \\
First-order condition & Assumed & Holds for $\mathrm{SU}(3)\times\mathrm{U}(1)_{B-L}$, fails for $\mathrm{SU}(2)_{L/R}$ \\
Coupling $\lambda$ & Not addressed & Free parameter, non-perturbative origin (\S\ref{sec:lambda}) \\
\bottomrule
\end{tabular}
\caption{Comparison with CCM noncommutative geometry approach~\cite{CCM2006}.}
\label{tab:ccm}
\end{table}

Three CCM postulates become derived results in this framework:
(1)~$\mathcal{A}_F=\CC\oplus\HH\oplus M_3(\CC)$ follows from $G_2$-harmonic analysis of $S^6$;
(2)~$N_{\mathrm{gen}}=3$ follows from the triality-channel index;
(3)~chirality follows from orientation of $S^6$ and $\mathrm{sign}(c_3)$.

\subsection{Comparison to Dolan--Nash \cite{DolanNash2002}}

Dolan and Nash derive Standard-Model fermion representations from
$\mathbb{CP}^3=\mathrm{SU}(4)/(\mathrm{SU}(3)\times\mathrm{U}(1))$ using CSDR. The two
constructions differ: internal space ($\mathbb{CP}^3$ vs.\ $S^6=G_2/\mathrm{SU}(3)$), generation
mechanism (Dolan-Nash derive SM representations but do not count generations from the Dirac
index), gauge group (Dolan-Nash require a 12-dimensional gauge group; present work uses spin
connection), and chirality (derived here from index sign, not addressed in Dolan-Nash). The
two constructions are complementary.

\subsection{Comparison to algebraic triality approaches \cite{FureyHughes2024,GillardGresnigt2019}}

Two other recent lines of work also derive three fermion generations from an $\mathrm{SO}(8)$
or $\mathrm{Spin}(8)$ triality structure, but through purely algebraic constructions rather than
a geometric index theorem on a continuous coset space. Furey and Hughes~\cite{FureyHughes2024}
identify the Standard Model within the triality symmetries of the division algebras
$\CC\oplus\HH\oplus\OC$; two generations arise directly as irreducible representations of a
``triality triple'', while the third emerges from a distinct ``Cartan factorization'' mechanism
in which division-algebraic multiplication merges spinor representations into a set of scalar
bosons. Gillard and Gresnigt~\cite{GillardGresnigt2019} construct three generations from the
$\mathrm{Cl}(8)$ Clifford algebra, with the three generations distinguished as inequivalent
complex structures on the octonions.

Both constructions are algebraic (division algebras / Clifford algebra representation theory)
and treat all three generations on an equal algebraic footing by construction. The present work
differs in deriving the three-channel structure from geometric data on a continuous
manifold --- the Atiyah--Singer index of the twisted Dirac operator on $S^6=G_2/\mathrm{SU}(3)$
--- and in leaving channel independence as an explicit open problem (L3b, Theorem~\ref{thm:elb3})
rather than resolving it algebraically. The two approaches are complementary rather than
competing: whether the $\mathrm{Spin}(8)$ fibre-symmetry input required by
Theorem~\ref{thm:elb3} admits an algebraic realization along the lines of
\cite{FureyHughes2024,GillardGresnigt2019} is an interesting open question, not addressed here.

Notably, the same obstruction recurs independently across formalisms. McRae~\cite{McRae2025}
identifies, at the level of $\mathrm{SO}(8)$ representation theory, that any triality-based
account of three generations must act at the level of representations rather than merely on
the Lie algebra, and discusses this as a standing obstacle rather than a solved problem.
Gourlay and Gresnigt~\cite{GourlayGresnigt2024} similarly work with an $S_3$ family symmetry
identified with the outer-automorphism triality of $\mathrm{Spin}(8)$, without a derivation of
this identification from an independent geometric or dynamical principle. Theorem~\ref{thm:elb3}
gives one concrete instance of why this is hard: on the specific coset geometry $G_2/\mathrm{SU}(3)$,
the three triality channels are provably indistinguishable to any $G_2$-equivariant construction.
This is offered as a parallel, independently obtained obstruction rather than a resolution of
these algebraic programs --- no bridge between the coset-bundle language used here and the
division-algebra or Clifford-algebra language used in
\cite{FureyHughes2024,GillardGresnigt2019,McRae2025,GourlayGresnigt2024} is established in this
work.

\subsection{Open problems}
\label{sec:open}

\textbf{Mathematical:}

\begin{enumerate}
\item \textbf{L3a --- $\ind=1$ per triality channel (\emph{proved}, this work).}
  All three SO(8) triality representations $\mathbf{8}_v, \mathbf{8}_s, \mathbf{8}_c$
  restrict to the same $G_2$-module $(1,0)_{G_2}\oplus(0,0)_{G_2}=7\oplus 1$ under
  $G_2=\mathrm{Fix}(\ZZ_3\subset\mathrm{Aut}(\mathrm{SO}(8)))$.
  Their $S^-$-subbundles have the same $\mathrm{SU}(3)$-module $(1,0)\oplus(0,0)$,
  so the Kostant--Parthasarathy argument of \S\ref{sec:schur} applies identically
  to every channel, giving $\ind(D^{(\alpha)})=1$ for $\alpha=0,1,2$
  (experiment E-L3-PARTIAL).

\item \textbf{L3b --- Channel independence (\emph{open; geometric path ruled out}).}
  Theorem~\ref{thm:elb3} (\S\ref{sec:bundle-obstruction}) proves that the $G_2$-equivariant
  path to channel independence is ruled out: $E_v\cong E_s\cong E_c$ as $G_2$-equivariant
  bundles, and no $G_2$-invariant operation --- including nearly-K\"ahler instanton
  connections --- can distinguish the three channels. The distinction requires
  $\mathrm{Spin}(8)$ fibre symmetry as an additional physical input. If the compactification
  setup preserves $\mathrm{Spin}(8)$ as an internal fibre symmetry, the three zero modes are
  $\mathrm{Spin}(8)$-invariantly orthogonal and $N_{\mathrm{gen}}=3$ follows. Confirming this
  physical input is the remaining open question, to be addressed in collaboration with
  T.~Lawrence.

\item \textbf{L4A --- Spectral computation of $\dim\ker(D\otimes S^-)=1$.}
  As argued in \S\ref{sec:lichnerowicz}, the Lichnerowicz bound does not directly give
  $\dim\ker=1$ (a uniform positive lower bound on $R/4+F_{S^-}$ would give $\ker=0$,
  contradicting $\ind=1$). The correct approach uses the Kostant--Parthasarathy formula
  for $(D\otimes S^-)^2$ on the naturally reductive space $G_2/\mathrm{SU}(3)$
  \cite{Agricola2002}, which expresses the squared Dirac eigenvalues in terms of $G_2$
  and $\mathrm{SU}(3)$ Casimir operators. This formula should yield $\dim\ker=1$ as a
  consequence of the $G_2$ representation theory.
  An independent route is the Casimir operator method of \cite{Dolan2003}, which computes
  Dirac spectra on any compact coset $G/H$ via quadratic Casimirs; it has been applied to
  $S^2, S^4, \mathbb{CP}^2$ and $\mathrm{SU}(3)/U(1)^2$ but not yet to
  $G_2/\mathrm{SU}(3)=S^6$.

\item \textbf{L4B --- $G_2$-equivariant character computation (\emph{partially resolved; rank
  on the trivial component supported by explicit computation, pending independent
  confirmation; non-trivial-sector vanishing for the Levi-Civita operator now directly
  verified for the two smallest representations, $\rho=\mathbf{7}$ established and
  $\rho=\mathbf{14}$ strongly supported}).}
  The Kostant--Parthasarathy formula (\S\ref{sec:schur}) proves a positive spectral gap on
  all non-trivial $G_2$-isotypic components of $\Gamma(S^+\otimes S^-)$ \emph{for Kostant's
  cubic Dirac operator}. For the physically relevant Levi-Civita operator, an independent,
  first-principles construction of the operator on matrix-coefficient sections (bypassing the
  cubic-operator framework entirely, diagonalizing its square on the full multiplicity space of
  each component) now directly confirms invertibility for $\rho=\mathbf{7}$ --- the smallest
  non-trivial Casimir gap, and the case singled out in \S\ref{sec:schur} as most exposed to a
  torsion correction --- and, with one flagged and honestly-scoped technical caveat (a relative
  sign not pinned down by any available internal consistency check, supported by a structural
  rationality argument but not an independent proof), for $\rho=\mathbf{14}$, the next-smallest
  gap. See \S\ref{sec:schur} for the full statement and its precise scope. On the trivial
  component, the map $D^+|_{\mathbf 1}:\CC^2\to\CC^1$ has index $1$ regardless of its rank
  (index of any linear map between spaces of these dimensions is automatically $2-1=1$), so the
  global index alone does not fix $\dim\ker(D^+_{S^-})\in\{1,2\}$. An explicit Clifford-algebra
  computation (Levi-Civita connection data of~\cite{AgrHofLawn2023}, calibrated against their
  Theorem~5.1 Killing-spinor eigenvalue, then applied to the twisted operator) gives rank$=1$
  ($\dim\ker(D^+_{S^-})=1$) directly, without relying on the naive Casimir formula that fails
  on this naturally reductive (non-symmetric) coset. This computation has passed calibration
  and two independent stress tests but has not yet received a final independent sign-off, so
  we report it as strong supporting evidence rather than a proved fact. See experiment
  \texttt{20260708-dolan-casimir-g2su3} (trivial-component rank computed; non-trivial-component
  torsion correction directly verified for $\rho=\mathbf{7},\mathbf{14}$; remaining
  representations formally open but structurally expected to be safe, per the growing-Casimir-
  gap argument of \S\ref{sec:schur}). The representation-theoretic setup builds
  on~\cite{AgrHofLawn2023,Agricola2002}.

\item \textbf{Integrability of $J$.} Lemma L1 uses the almost-complex structure on $S^6$, which
  is non-integrable (Borel--Serre theorem). The Atiyah--Singer theorem applies to almost-Hermitian
  manifolds, but curvature contributions to the Weitzenb\"ock formula may differ from the integrable
  case. This must be verified for the nearly K\"ahler setting of $S^6$.

\end{enumerate}

\textbf{Physical:}

\begin{enumerate}
\setcounter{enumi}{3}
\item \textbf{$\rho_3$ modulus stabilization.} The potential $V_{\mathrm{total}}$ depends on $\rho_3$
  only through $C^3$, creating no restoring force in the $\rho_3$ direction (Dine-Seiberg runaway).
  Stabilizing $\rho_3$ requires an additional contribution (D-term, flux threading a 3-cycle of $S^3$,
  or brane action).

\item \textbf{Non-perturbative origin of $\lambda$.} As proved in \S\ref{sec:lambda}, no geometric
  mechanism produces $e^{-\lambda/\rho_6^2}$. A complete moduli stabilization requires specifying
  the string background and computing $\lambda$ from first principles.

\item \textbf{Fermion mass hierarchy.} The index theorem gives $N_{\mathrm{gen}}=3$ with equal
  geometric status for each generation; it does not explain why $m_t/m_u\approx 10^5$.
  In fact, a stronger statement holds: Theorem~\ref{thm:yukawa-deg} (\S\ref{sec:yukawa-deg})
  shows --- conditional on the L4B kernel-rank assumption (\S\ref{sec:kernel}) --- that all
  inter-generation Yukawa couplings are equal on $S^6$, so the mass
  hierarchy must originate from the non-geometric $D_F$ freedom of the finite spectral
  triple \cite{CCM2006}. The role of $S^3$ Kaluza-Klein coupled-mode mixing in lifting
  the degeneracy in the full 9-dimensional analysis remains an open question.

\item \textbf{Universality.} Whether the Lichnerowicz--$G_2$-Schur mechanism (L4) applies to other
  nearly K\"ahler 6-manifolds (such as $\mathbb{CP}^3$ or $\mathrm{SU}(3)/T^2$) is an open question.

\item \textbf{Strong CP problem.} The $S^3$ Dirac operator has spectrum $\{\pm(n+3/2): n\geq 0\}$
  --- perfectly paired positive and negative eigenvalues. The Atiyah--Patodi--Singer $\eta$-invariant
  therefore vanishes: $\eta(0, D_{S^3})=0$. By the APS index theorem, a vanishing $\eta$-invariant
  implies no CP-odd boundary contribution from the $S^3$ sector to the topological QCD theta angle.
  Whether this extends to a complete resolution of $\theta_{QCD}=0$ requires a global Pontryagin
  density analysis on $S^3\times S^6$ including $S^6$ sector and non-perturbative effects.

\item \textbf{$\lambda$ as topological ratio [HYPOTHESIS].}
  The free coupling satisfies $\lambda = N_{\mathrm{gen}}/\dim(S^3\times S^6) = 3/9 = 1/3$
  if $\lambda$ is identified with the ratio of the generation count to the total internal dimension.
  Unlike geometric derivations from the radii $\rho_3, \rho_6$ (forbidden by Corollary C1), a
  topological identification of this kind does not require $\lambda$ to be expressible as
  $c\cdot\rho_3^a\cdot\rho_6^{2-a}$. Whether any known brane instanton or gaugino condensation
  mechanism produces exactly $\lambda=1/3$ is unknown and left for future work.
\end{enumerate}

%% ================================================================
\section*{Acknowledgements}

The author thanks Tom Lawrence for numerous discussions on the $S^3$ Kaluza-Klein mechanism
and for the foundational framework of~\cite{Lawrence2022}.

The author acknowledges the use of Claude (Anthropic) as a computational and drafting
assistant throughout this work, including symbolic and numerical verification of the
derivations (implemented as a Python/\texttt{sympy} test suite, 2484 tests, all passing),
literature cross-checks, and manuscript drafting assistance. All mathematical claims,
proofs, and numerical results were independently reviewed and verified by the author,
who takes full responsibility for the content of this paper, consistent with arXiv's
policy on the use of generative AI tools.

\medskip
\noindent\textbf{Code availability.} The complete verification suite accompanying this
work (exact-arithmetic \texttt{sympy} derivations and regression tests for every claim)
is available at \url{https://github.com/sergeeey/N-7-GeoSpectra-Lab}
(directory \texttt{tom\_s3\_spinor\_toy}).

%% ================================================================
\bibliographystyle{unsrt}
\begin{thebibliography}{99}

\bibitem{Lawrence2022}
T.~Lawrence,
\textit{Product manifolds as realisations of general linear symmetries},
arXiv:2203.09473 (2022).

\bibitem{DolanNash2002}
B.P.~Dolan and C.~Nash,
\textit{Chiral fermions from spinors of SO(8)},
JHEP \textbf{10} (2002) 041.

\bibitem{FureyHughes2024}
N.~Furey and M.J.~Hughes,
\textit{Three generations and a trio of trialities},
Phys.\ Lett.\ B (2025), arXiv:2409.17948.

\bibitem{GillardGresnigt2019}
A.B.~Gillard and N.G.~Gresnigt,
\textit{The $\mathrm{Cl}(8)$ algebra of three fermion generations with spin and full internal symmetries},
arXiv:1906.05102 (2019).

\bibitem{McRae2025}
C.~McRae,
\textit{Exploring triality explicitly: Convenient bases for $\mathrm{SO}(8)$, $\mathrm{Spin}(1,7)$, and $G_2$},
arXiv:2502.14016 (2025).

\bibitem{GourlayGresnigt2024}
L.~Gourlay and N.G.~Gresnigt,
\textit{Algebraic realisation of three fermion generations with $S_3$ family and unbroken gauge symmetry from $\mathrm{C}\ell(8)$},
Eur.\ Phys.\ J.\ C \textbf{84}, 1129 (2024), arXiv:2407.01580.

\bibitem{CCM2006}
A.H.~Chamseddine, A.~Connes, and M.~Marcolli,
\textit{Gravity and the Standard Model with neutrino mixing},
arXiv:hep-th/0610241 (2006).

\bibitem{ChamseddineConnes1997}
A.H.~Chamseddine and A.~Connes,
\textit{The spectral action principle},
Commun.\ Math.\ Phys.\ \textbf{186} (1997) 731--750.

\bibitem{HarlandNolle2011}
D.~Harland and C.~N\"olle,
\textit{Instantons and Killing spinors},
arXiv:1109.3552 (2011).

\bibitem{AtiyahSinger1963}
M.F.~Atiyah and I.M.~Singer,
\textit{The index of elliptic operators on compact manifolds},
Bull.\ Amer.\ Math.\ Soc.\ \textbf{69} (1963) 422--433.

\bibitem{Milnor1963}
J.~Milnor,
\textit{Morse Theory},
Princeton University Press, 1963.

\bibitem{Buckingham1914}
E.~Buckingham,
\textit{On physically similar systems; illustrations of the use of dimensional equations},
Phys.\ Rev.\ \textbf{4} (1914) 345--376.

\bibitem{Butruille2005}
J.-B.~Butruille,
\textit{Classification des vari\'et\'es approximativement k\"ahl\'eriennes homog\`enes},
Ann.\ Global Anal.\ Geom.\ \textbf{27} (2005), 201--225;
arXiv:math/0401152.

\bibitem{KN1969}
S.~Kobayashi and K.~Nomizu,
\textit{Foundations of Differential Geometry}, Vol.~II,
Wiley Interscience (1969), \S 8.
(For any almost-Hermitian manifold $(M,g,J)$: $c_n(T^{1,0}M)=e(TM)$,
integrability of $J$ not required.)

\bibitem{AgrHofLawn2023}
I.~Agricola, T.~Hofmann, and M.-A.~Lawn,
\textit{Invariant spinors on homogeneous spheres},
Differential Geom.\ Appl.\ \textbf{88} (2023), 102014;
arXiv:2203.02961.
(Theorem~5.1: $S^6=G_2/\mathrm{SU}(3)$ has exactly two $G_2$-invariant Killing spinors
$\psi_\pm = 1\pm y_1\wedge y_2\wedge y_3$.)

\bibitem{Agricola2002}
I.~Agricola,
\textit{Connections on naturally reductive spaces, their Dirac operator and homogeneous
models in string theory},
Commun.\ Math.\ Phys.\ \textbf{232} (2002) 535--563;
arXiv:math/0202094.
(Theorems~3.3 and~4.2: Kostant--Parthasarathy formula for $(D^t)^2$ on $G/H$;
$G$-invariant spinors are eigenspinors with eigenvalue $9t^2[Q(\rho_g,\rho_g)-Q(\rho_h,\rho_h)]>0$.)

\bibitem{Dolan2003}
B.P.~Dolan,
\textit{The Spectrum of the Dirac Operator on Coset Spaces with Homogeneous Gauge Fields},
arXiv:hep-th/0304037 (2003).
(Casimir operator method for Dirac spectra on compact $G/H$; explicit results for
$S^2, S^4, \mathbb{CP}^2, \mathrm{SU}(3)/U(1)^2$; method applies to any compact coset.)

\end{thebibliography}

\end{document}
